% =============================================================================
% Chapter 9: Contributing Guide
% =============================================================================

\chapter{Contributing Guide}
\label{ch:contributing}

\section{Overview}

\itc{} handles data that feeds engineering decisions. Every contribution
must meet a higher standard than typical open-source software. This
chapter defines that standard. First-time contributors should read
Chapters~\ref{ch:philosophy} and~\ref{ch:datamodel} before writing any
code.

\section{Getting Started}

\subsection{Repository Setup}

\begin{lstlisting}[language=bash]
# Clone the repository
git clone https://github.com/nevesgeovana/itaca.git
cd itaca

# Create a virtual environment
python -m venv .venv
source .venv/bin/activate          # Linux / macOS
# .venv\Scripts\activate           # Windows

# Install in editable mode with all development dependencies
pip install -e ".[dev]"

# Install pre-commit hooks
pre-commit install

# Verify the setup
pytest --cov=itaca --cov-report=term-missing
ruff check itaca/
ruff format --check itaca/
mypy --strict itaca/
\end{lstlisting}

\subsection{Development Dependencies}

The \code{[dev]} extra installs everything in
Table~\ref{tab:dependency-versions} marked as \code{dev}, plus the full
set of optional runtime dependencies needed to test the optional code
paths.

\section{Workflow}

\subsection{Branch Strategy}

\begin{itemize}[noitemsep,topsep=2pt]
  \item \code{main}: stable, always passing CI, never commit directly.
  \item \code{dev}: integration branch for completed features.
  \item \code{feat/short-description}: feature branches, branched from \code{dev}.
  \item \code{fix/short-description}: bug-fix branches.
  \item \code{docs/short-description}: documentation-only changes.
\end{itemize}

\subsection{The TDD Loop}

Every new public method or feature follows this exact order:

\begin{enumerate}
  \item \textbf{Write the usage example.} Before writing any code, write
    how the feature will be called. This becomes both the docstring
    example and an integration test.
  \item \textbf{Write the test.} Write \code{pytest} tests that specify
    the expected behavior, including the relevant edge cases
    (Table~\ref{tab:required-edge-cases}). Where mathematical contracts
    apply, write a property-based test using
    Hypothesis~\citep{hypothesis}. The tests must fail at this point.
  \item \textbf{Write the implementation.} Write the minimum code needed
    to make the tests pass.
  \item \textbf{Refactor.} Clean up the implementation without breaking
    the tests.
  \item \textbf{Verify coverage.} Run \code{pytest --cov=itaca} and
    confirm coverage is $\ge 90\%$ on changed code.
\end{enumerate}

\paragraph{Worked example: adding \code{db.squeeze()}.}

The following walks through the full loop for a small but real
operation. It exists in \itc{} as REQ-22; the example reconstructs the
discipline that produced it.

\textbf{Step 1: usage example, written first.}
\begin{lstlisting}
db = itc.load("data/", dims=["alpha", "mach", "blade"])
db_at = db.at(mach=0.12)            # mach is now cardinality-1
print(db_at)                         # alpha(N) x mach(1) x blade(M)

db_clean = db_at.squeeze()
print(db_clean)                      # alpha(N) x blade(M)

# squeeze a specific dim only
db_explicit = db_at.squeeze(along="mach")
\end{lstlisting}

\textbf{Step 2: tests written before implementation.}
\begin{lstlisting}
# tests/ops/test_squeeze.py
import numpy as np
import pytest
import itaca as itc
from itaca.core.errors import DimensionNotFoundError


def test_squeeze_removes_unit_dim(simple_varframe_unit_mach):
    db = simple_varframe_unit_mach
    out = db.squeeze()
    assert "mach" not in out.dims
    assert set(out.dims) == set(db.dims) - {"mach"}


def test_squeeze_preserves_history(simple_varframe_unit_mach):
    db = simple_varframe_unit_mach
    out = db.squeeze(comment="drop singleton mach")
    last = out.history.entries[-1]
    assert last.operation.startswith("squeeze")
    assert last.comment == "drop singleton mach"


def test_squeeze_along_specified_dim(simple_varframe_unit_mach):
    db = simple_varframe_unit_mach
    out = db.squeeze(along="mach")
    assert "mach" not in out.dims


def test_squeeze_along_non_unit_raises(simple_varframe):
    db = simple_varframe
    with pytest.raises(ValueError):
        db.squeeze(along="alpha")     # alpha cardinality > 1


def test_squeeze_unknown_dim_raises(simple_varframe):
    with pytest.raises(DimensionNotFoundError):
        simple_varframe.squeeze(along="not_a_dim")


def test_squeeze_full_collapse_yields_datapoint(scalar_varframe):
    out = scalar_varframe.squeeze()
    assert list(out.dims) == ["datapoint"]
    assert out.dims["datapoint"].coords.size == 1
\end{lstlisting}

\textbf{Step 3: minimal implementation.}
\begin{lstlisting}
# itaca/ops/squeeze.py
from __future__ import annotations
import numpy as np
from itaca.core.errors import DimensionNotFoundError


def squeeze(self, along: str | None = None,
            comment: str | None = None) -> "VarFrame":
    if along is not None and along not in self.dims:
        raise DimensionNotFoundError(
            f"Dimension '{along}' not found in VarFrame.\n"
            f"Operation: squeeze(along='{along}')\n"
            f"Available: {list(self.dims)}\n"
            f"Fix: pass an existing dimension or omit `along`."
        )

    targets = [along] if along else [
        d for d, dim in self.dims.items() if dim.coords.size == 1
    ]
    for d in targets:
        if self.dims[d].coords.size != 1:
            raise ValueError(
                f"Dimension '{d}' has cardinality "
                f"{self.dims[d].coords.size}, cannot squeeze."
            )

    new_dims  = {d: dim for d, dim in self.dims.items() if d not in targets}
    new_vars  = {n: v.squeeze(axis=tuple(self._axis(d) for d in targets))
                 for n, v in self.vars.items()}
    new_tags  = self.tags.squeeze(targets) if self.tags else None
    new_unc   = self.uncertainty.squeeze(targets) if self.uncertainty else None
    new_hist  = self.history.append(
        operation=f"squeeze(along={targets!r})", comment=comment,
    )
    if not new_dims:
        new_dims = self._datapoint_dim(size=1)

    return self._replace(dims=new_dims, vars=new_vars, tags=new_tags,
                         uncertainty=new_unc, history=new_hist)
\end{lstlisting}

\textbf{Step 4: refactor.} Extract the axis-resolution helper into
\code{core/index.py} so \code{average} and \code{integrate} can reuse it.
Tests still pass.

\textbf{Step 5: coverage.} \code{pytest --cov=itaca.ops.squeeze}
reports 100\,\%; the property-based test in
\code{test\_ops\_properties.py} confirms that
\code{db.expand(d, [v]).squeeze(along=d)} round-trips for any
randomly-generated VarFrame.

\begin{noteinfo}[Pull requests without tests are rejected]
This is a hard rule, regardless of correctness. The TDD loop is the
project's only mechanism for guaranteeing that the SRS and the code stay
aligned.
\end{noteinfo}

\section{Coding Standards}

\subsection{File and Module Structure}

Each module has a single, clearly stated responsibility described in its
module docstring. Modules longer than 400 lines should be considered for
splitting.

\begin{lstlisting}
"""
itaca.core.varframe
===================

Defines the VarFrame class: the primary multidimensional data
structure in ITACA.

All operations on VarFrame are immutable: they return new objects and
never modify the original. Provenance is preserved; History is appended.
"""
from __future__ import annotations
import numpy as np
from itaca.core.provenance import Provenance
from itaca.core.history import History
\end{lstlisting}

\subsection{Type Hints and Docstrings}

All public functions and classes carry full type hints and NumPy-format
docstrings. The following is the mandatory template:

\begin{lstlisting}
def select(self, filters: dict, *,
           Frame: str = "VarFrame",
           comment: str | None = None) -> "VarFrame":
    """
    Return a VarFrame restricted to the dimension values matching the filters.

    Parameters
    ----------
    filters : dict
        Mapping of ``dim_op -> value`` where ``dim_op`` is the dimension
        name optionally suffixed by an operator (``=``, ``>``, ``<``,
        ``>=``, ``<=``, ``!=``). Default operator is ``=``.
    Frame : {"VarFrame", "HistoryFrame", "UncFrame"}, optional
        Source against which the comparison is evaluated. Default
        compares against the variable values.
    comment : str, optional
        Optional comment recorded in the History entry.

    Returns
    -------
    VarFrame
        New VarFrame containing only the specified coordinate values.
        Provenance is preserved; History gains one entry for this operation.

    Raises
    ------
    DimensionNotFoundError
        If any filter key does not match a dimension name.
    SelectionError
        If any specified value does not exist in the dimension coordinates.

    Examples
    --------
    >>> db_sub = db.select({"alpha": [0, 2, 4], "mach>": 0.15})
    >>> db_orig = db.select({"CT": 0}, Frame="HistoryFrame")
    """
\end{lstlisting}

\subsection{Error Messages}

Every raised exception includes three pieces of information:

\begin{lstlisting}
# Bad: terse, unhelpful
raise KeyError(f"Variable not found: {name}")

# Good: variable name, operation, fix suggestion
raise VariableNotFoundError(
    f"Variable '{name}' not found in VarFrame.\n"
    f"Operation: compute('{name} = ...')\n"
    f"Available variables: {list(self.vars.keys())}\n"
    f"Fix: check the variable name or load the source data containing it."
)
\end{lstlisting}

All \itc{}-specific exceptions inherit from \code{ITACAError}
(Table~\ref{tab:exception-hierarchy}).

\subsection{History Recording}

Every method that returns a new VarFrame appends an entry to the History
and updates the HistoryFrame when applicable:

\begin{lstlisting}
def fill(self, along: str, method: str = "linear",
         comment: str | None = None, **kwargs) -> "VarFrame":
    # ... computation ...
    new_history = self.history.append(
        operation=f"fill(along='{along}', method='{method}')",
        kwargs=kwargs,
        comment=comment,
    )
    new_tags = self.tags.mark_filled(filled_mask)        # set +1 where filled
    return self._replace(vars=new_vars, tags=new_tags, history=new_history)
\end{lstlisting}

\subsection{NumPy-Only in Core}

The \code{core/}, \code{ops/}, and \code{uncertainty/} packages import
only from the Python standard library and NumPy. A \code{ruff}
import-policy rule flags any other import in these packages.

\subsection{Language and Orthography}

All code, comments, docstrings, variable names, error messages, and
\code{.itceq} files are in English. American English with Z is used
consistently (\emph{organize}, \emph{visualization}, \emph{normalize}).

\subsection{Example and Test Data}

Every dataset that ships with the repository (examples, test fixtures,
documentation) is either synthetic or publicly licensed, and its
provenance is stated where it is used. Employer-origin, proprietary, or
unreleased experimental data never enter the repository, in any form.
Contributors publishing artifacts generated with \itc{} should also note
that exports embed the user identity from \code{getpass} and the machine
hostname by default; \code{itc.set\_user(...)} sets a publication-safe
identity.

\section{Commit Messages}

Follow Conventional Commits:

\begin{tabularx}{\textwidth}{@{}lX@{}}
\toprule
\textbf{Prefix} & \textbf{Use for}\\
\midrule
\code{feat:}     & New feature or public API addition\\
\code{fix:}      & Bug fix\\
\code{docs:}     & Documentation only\\
\code{test:}     & Test addition or modification\\
\code{refactor:} & Code restructuring without behavior change\\
\code{perf:}     & Performance improvement\\
\code{build:}    & Build system or external dependencies\\
\code{ci:}       & CI configuration\\
\code{chore:}    & Maintenance tasks (dependency bumps, etc.)\\
\bottomrule
\end{tabularx}

\begin{lstlisting}[language=bash]
feat: add fitmodel() returning polynomial coefficients as new dim
fix: correct polar integration area element in integrate()
test: add property-based test for RPN sin/cos derivative identities
docs: update contributing guide with HistoryFrame update pattern
refactor: split varframe.py into varframe.py and dimension.py
\end{lstlisting}

The commit subject line shall not exceed 72 characters. The body
explains \emph{why}, not what (the diff already shows what).

\section{Pull Request Checklist}

Before opening a pull request, every item in this checklist must be
verified:

\begin{itemize}[leftmargin=2.5em,label=$\square$]
  \item Usage example written and added to docstring
  \item Tests written before implementation
  \item All new tests pass: \code{pytest}
  \item Property-based test added if a math contract applies
  \item Coverage $\ge 90\%$: \code{pytest --cov=itaca}
  \item Linting passes: \code{ruff check itaca/}
  \item Formatting correct: \code{ruff format --check itaca/}
  \item Type-checking passes: \code{mypy --strict itaca/}
  \item Pre-commit hooks pass: \code{pre-commit run --all-files}
  \item Type hints present on all public signatures
  \item Docstring in NumPy format with Examples section
  \item Every new VarFrame-returning method records History and updates
        the HistoryFrame when applicable
  \item No pandas/xarray/dask imports in \code{core/}, \code{ops/}, or
        \code{uncertainty/}
  \item All code, comments, and docstrings in American English with Z
  \item Example and test data synthetic or publicly licensed, provenance stated
  \item \code{CHANGELOG.md} updated if public API changed
  \item Commit messages follow Conventional Commits
\end{itemize}

\section{Adding a New Built-in Processor}

The following walks through adding a hypothetical
\code{ft\_drag\_polar} processor, end to end. It is structured to mirror
the TDD loop above.

\paragraph{Step 1: usage example.}
\begin{lstlisting}
pproc = itc.pproc("ft_drag_polar", config={"S_ref": 50.0})
pproc.info()                         # prints: equations, uncertainties, required vars
pproc.validate(db)                   # raises ProcessorValidationError if vars missing
db_out = pproc(db, comment="initial polar from FT-2026-A")
db_out = pproc(db, report="reports/ft_drag.pdf")
\end{lstlisting}

\paragraph{Step 2: companion \code{.itceq}.} Create
\code{itaca/pproc/equations/flight\_test/ft\_drag\_polar.itceq} with the
five sections (Section~\ref{sec:itceq}). Keep \code{[meta]} accurate and
list every measured-variable uncertainty in
\code{[uncertainties]}. Equations evaluate in file order by default;
enable \code{auto\_sort=True} only if needed.

\paragraph{Step 3: processor implementation.}
\begin{lstlisting}
# itaca/pproc/builtin/ft_drag_polar.py
from __future__ import annotations
from itaca.pproc.base import EquationProcessor

class FTDragPolar(EquationProcessor):
    name        = "ft_drag_polar"
    version     = "0.1.0"
    idempotent  = False

    def __init__(self, config: dict | None = None):
        super().__init__(
            equations_file="flight_test/ft_drag_polar.itceq",
            config=config or {},
        )

    def info(self) -> None:
        # Prints the parsed metadata, constants, and equations.
        super().info()

    def validate(self, db) -> None:
        super().validate(db)
        # Additional domain-specific checks go here.
\end{lstlisting}

\paragraph{Step 4: registry entry.} Append to
\code{itaca/pproc/registry.py}:
\begin{lstlisting}
from itaca.pproc.builtin.ft_drag_polar import FTDragPolar
PROCESSORS["ft_drag_polar"] = FTDragPolar
\end{lstlisting}

\paragraph{Step 5: tests.} Create
\code{tests/pproc/test\_ft\_drag\_polar.py} covering: happy path on a
synthetic VarFrame fixture; missing required variable raises
\code{ProcessorValidationError}; custom config overrides defaults;
History records exactly one entry per application; report generation
produces a PDF; idempotence guard fires on the second application unless
\code{force=True}.

\paragraph{Step 6: documentation.} Add a short worked example to
\code{docs/processors.md} and update the table in
Section~\ref{ch:functional} listing built-in processors.

\section{Adding a Specialized Plot Method}

\begin{enumerate}[noitemsep,topsep=2pt]
  \item Create \code{itaca/plot/specialized/my\_plot.py}.
  \item Implement the method as a function accepting a \code{DataVis}
    instance as \code{self} (mixin pattern).
  \item Register it as a method on \code{DataVis} in
    \code{itaca/plot/datavis.py}.
  \item Write tests using a real VarFrame fixture: never synthetic
    random data.
\end{enumerate}

\section{License and Citation}

\itc{} is released under the MIT License. The license file lives at the
repository root; every source file is covered by it without per-file
headers. Contributions are accepted under the same license.

Tagged releases are mirrored on Zenodo with a DOI (REQ-97). The
repository root carries a \code{CITATION.cff} file so that GitHub and
Zenodo expose a canonical citation; the file is updated at every release
together with \code{CHANGELOG.md}. After the first stable release, a
software paper (JOSS or SoftwareX) is planned as the citable reference
(Chapter~\ref{ch:roadmap}).

\section{SRS Document Maintenance}

This SRS is version-controlled alongside the code. When adding
requirements:

\begin{enumerate}[noitemsep,topsep=2pt]
  \item Assign the next available \reqid{REQ-XX} number (do not reuse
    deprecated identifiers).
  \item Set status to \draft{} until the corresponding code is
    implemented and tested.
  \item Update the revision history table on the title page.
  \item Increment the document version following semantic versioning
    (\code{MINOR} for new requirements, \code{PATCH} for clarifications).
\end{enumerate}

Never delete a requirement: mark it \deprecated{} and explain why.
